% Executable format reference
% Keep this chapter grounded in the loader and assembler implementations.
\chapter{Executable File Format}
\label{ch:executable-format}

\todo{Specify the executable format from the actual loader and assembler implementations. Do not infer header fields, load addresses, relocation behavior, or entry-point rules from examples alone.}

This chapter will describe the on-disk representation of executable programs
and the rules used by the system loader. The specification must be checked
against the implementation before it is treated as authoritative.

\section{Format overview}
\todo{Document magic bytes or signatures, header layout, versioning, and how the loader distinguishes executable files from other file types.}

\section{Header and payload}
\todo{Provide a byte-offset table for every header field, including size, encoding, valid ranges, and interpretation. Document payload placement and any alignment requirements.}

\section{Loading and execution}
\todo{Trace the loader's validation, memory allocation, load-address selection, relocation (if any), entry-point selection, and error paths.}

\section{Producing executable files}
\todo{Document the assembler or other toolchain output options, with reproducible examples and a byte-for-byte annotated sample.}

\section{Compatibility and failure cases}
\todo{Record format versions, compatibility rules, malformed-file handling, and relevant limits.}
